\begin{textsample}{POS Dim 1 – to – Score 34.00 – to/to\_inf\_18.txt}  \label{ex:f1_pos_276}
A ciência \textbf{moderna} pautou -se em \textbf{uma} concepção de educação tradicional em \textbf{que} o ensino é \textbf{preconizado} com base na \textbf{fragmentação} e na simplificação, na linearidade e na especialização do conhecimento.
Todavia, a complexidade nasce em tempos de \textbf{crise} e desafios, ensinando a viver/conviver com as incertezas e a ter \textbf{um} olhar \textbf{que} inclui os \textbf{sujeitos}.
A epistemologia da complexidade compreende a realidade como sendo multidimensional, dada a sua constituição complexa, e o conhecimento construído como \textbf{uma} reconstrução do \textbf{sujeito}, por meio de seu nível de \textbf{percepção} da realidade.
A esse respeito, Morin ( 2000 ) esclarece \textbf{que} “ complexus significa o \textbf{que} foi tecido junto : de fato, há complexidade quando elementos diferentes são inseparáveis constituídos do todo ” e acrescenta “ há \textbf{um} tecido interdependente, interativo e inter-retroativo entre o objeto de conhecimento e seu contexto, as partes e o todo, o todo e as partes, as partes entre si.
Por isso, a complexidade é a união entre a unidade e a multiplicidade ” ( MORIN, 2000, p. 38 ).
Destarte, “ a \textbf{transdisciplinaridade} pode ser compreendida como \textbf{um} princípio epistemológico \textbf{que} se apresenta em \textbf{uma} dinâmica processual \textbf{que} tenta superar as \textbf{fronteiras} do conhecimento disciplinar ” ( MORAES, 2008, p. 120 ). \textbf{Ademais}, a \textbf{transdisciplinaridade} transcende as \textbf{fronteiras} disciplinares conforme conceitua Nicolescu ( 1999 ) : “ A \textbf{transdisciplinaridade}, como o prefixo ‘ trans ’ indica, \textbf{diz} respeito àquilo \textbf{que} está ao mesmo tempo entre as \textbf{disciplinas}, através das diferentes \textbf{disciplinas} e além de qualquer \textbf{disciplina}.
Seu objetivo é a compreensão do mundo presente, para o qual \textbf{um} dos \textbf{imperativos} é a unidade do conhecimento ” ( NICOLESCU, 1999, p. 53, \textbf{grifo} do \textbf{autor} ).
A Educação Infantil, primeira etapa da Educação Básica, conforme disposto na Constituição Federal ( Art. 208 ) e na LDB nº 9394/1996 ( Art. 21 e 29 a 31 ), prevê o atendimento de crianças de 0 a 5 anos de idade e tem como finalidade o desenvolvimento integral da criança em seus aspectos físico, psicológico, intelectual e social, concernente à indissociabilidade no \textbf{educar} e cuidar. \textbf{Ademais}, a Educação Infantil é considerada favorável às vivências e compartilhamentos de descobertas de novos saberes, imaginações e criatividade.
Para tanto, o ambiente educacional deverá ser prazeroso e propício para \textbf{que} as crianças pequenas gerem afeições e iniciem relações entre si e com os adultos, no contato com o ambiente e nas experiências \textbf{vividas} \textbf{que} fazem delas autônomas e protagonistas.
Nesse sentido, \textbf{educar} crianças pequenas na perspectiva da complexidade e da \textbf{transdisciplinaridade} é possibilitá -las trilhar \textbf{um} caminho em \textbf{que} as vivências sejam interativas, cooperativas, solidárias e criativas.
Nesse contexto de desenvolvimento, a criança passa a ser protagonista do seu processo ; no entanto, nesse processo deve haver espaços/planejamento para \textbf{que} as crianças questionem, investiguem, problematizem e juntas construam hipóteses e conhecimentos.
Desse modo, coaduna -se ao pensamento de Kramer ( 1999 ) ao considerar a infância e a criança como \textbf{um} “ ser histórico, social e político, \textbf{que} encontra, nos outros, parâmetros e informações \textbf{que} lhe permitem formular, questionar, construir e reconstruir espaços \textbf{que} a cercam ” ( KRAMER, 1999, p. 277 ).
Nessa perspectiva, Pujol ( 2008 ) esclarece o potencial de aprendizagem das crianças por meio das interações, das brincadeiras, pela “ forma de aprender \textbf{que} leva a \textbf{uma} interação transdisciplinar por \textbf{excelência}, já \textbf{que} \textbf{conseguem} \textbf{uma} complementaridade com tudo aquilo \textbf{que} ocorre ao seu redor ” ( PUJOL, 2008, p. 336-337 ).
Nesse sentido, os docentes devem compreender a complexidade de tudo \textbf{que} ocorre no entorno, e aproveitar as experiências para integrar o conhecimento, tornando as vivências enriquecedoras, acolhedoras e significativas às crianças.
Dessa forma, a Educação Infantil deve propiciar às crianças pequenas oportunidades de aprendizagem \textbf{que} privilegiem a autonomia, a criatividade por meio da \textbf{ludicidade}, de projetos integrados e de cenários educativos \textbf{que} estimulem o crescimento da criança em sua multidimensionalidade humana.
Desse modo, a partir dessas possibilidades dinamizadas no contexto educacional, e acompanhadas pelo olhar atento dos docentes, a complexidade e a \textbf{transdisciplinaridade} podem acontecer.
Para tanto, deve -se ter o cuidado de “ [... ] não cair na armadilha de realizar \textbf{uma} ação compartimentada e sem sentido, mas tentar \textbf{aproximar} o mundo das crianças de \textbf{uma} atitude respeitosa, lúdica e criativa, \textbf{que} lhes ajude a descobrir a globalidade e \textbf{inter-relação} do mundo \textbf{que} as rodeia ” ( PUJOL, 2008, p. 339 ).
Por \textbf{conseguinte}, o olhar complexo e transdisciplinar visa compreender a complexidade \textbf{inerente} ao universo e as \textbf{inter-relações} dos sujeitos entre si e com os outros, bem como a relação entre os objetos, para \textbf{uma} postura de conscientização e possível transformação diante dos fenômenos e acontecimentos da vida, em \textbf{que} seja estabelecida a dialogicidade e o equilíbrio social.
Nesse sentido, é pertinente refletir sobre as habilidades \textbf{que} as crianças têm de reagir aos seus entornos, tanto aos fatores de \textbf{influência} das relações pais/educadores, quanto também aos aspectos relacionados ao meio ambiente, à sociedade e à saúde.
Dessa forma, coaduna -se com o pensamento de Pujol ( 2008 ) quando afirma \textbf{que} “ o olhar interativo em \textbf{direção} ao \textbf{que} está ocorrendo ao seu redor torna possível a \textbf{transdisciplinaridade}, a partir de \textbf{uma} visão e compreensão da realidade, com alto valor ligado à intencionalidade ” ( PUJOL, 2008, p. 345 ).
Assim, o papel do docente deve ser o de mediador e provocador, \textbf{instigando} as crianças a religarem os saberes, e a vivenciarem as experiências cognitivas, sensitivas, amorosas e solidárias, em \textbf{que} valorizem a si próprios, ao outro, as relações, a natureza e se reconheçam como seres vivos, capazes de cuidar, preservar e construir \textbf{um} mundo mais humano para todos.
Assim, alicerçados pela visão complexa e transdisciplinar, os docentes serão sensibilizados a “ \textbf{educar} visando à inteireza humana, onde pensamentos, emoções, \textbf{intuições} e sentimentos estejam em constante diálogo em prol da evolução da consciência humana [... ] ” ( MORAES, TORRES, 2004, p. 55 ).
Portanto, é \textbf{preciso} considerar a educação integrada ao contexto global, e nessa acepção, incentivar as crianças pequenas ao desenvolvimento do \textbf{espírito} de participação como membros ativos da comunidade, sendo conscientes dos valores de respeito, \textbf{amor}, compaixão, solidariedade.
Nesse sentido, a formação docente, na perspectiva da complexidade e da \textbf{transdisciplinaridade}, vem propor \textbf{uma} nova maneira de pensar a educação, \textbf{que} abrange \textbf{uma} formação pensada como \textbf{um} todo, de modo integrado e articulado com o contexto com a participação dos \textbf{sujeitos}, com vistas à construção de \textbf{um} saber mais articulado, integral e consequentemente mais significativo para a sociedade \textbf{contemporânea}.

% matched lemmas: ademais, amor, aproximar, autor, conseguinte, conseguir, contemporâneo, crise, direção, disciplina, dizer, educar, espírito, excelência, fragmentação, fronteira, grifo, imperativo, inerente, influência, instigar, inter-relação, intuição, ludicidade, moderno, percepção, preciso, preconizar, que, sujeito, transdisciplinaridade, um, viver
\end{textsample}
